\section{Amplificatori Operazionali Reali - Seconda Parte - 30.09.2014}

In questa esperienza abbiamo studiato ulteriori caratteristiche che distinguono l'opamp reale dal modello ideale.
Più precisamente abbiamo analizzato e misurato lo \textit{slew rate}, la massima corrente pilotabile, la larghezza di banda passante e il guadagno a maglia aperta $A_{ol}$ di un opamp $\mu$A741.

Ricordiamo che in ogni schema sono stati sottintesi la circuitazione di alimentazione (cavi e capacitori) e quella dedicata al bilanciamento dell'offset (cavi e trimmer).

\subsection*{Materiali}

\begin{itemize} [noitemsep]
%\item Oscilloscopio Agilent DSO-X 2002A (bandwidth \SI{70}{\mega\hertz}, sample rate \num{2} GSa/s);
%\item Generatore di tensione continua Agilent E3631A (max $\pm \, \SI{25}{\volt}$ o $\pm \, \SI{6}{\volt}$);
%\item Generatore di forme d'onta Agilent 33120A con range di frequenza da \SI{100}{\micro\hertz} a \SI{15}{\mega\hertz};
%\item Multimetro Agilent 34410A;
\item Un amplificatore operazionale $\mu$A741;
\item Resistenze e capacità di vari valori;
%\item un trimmer (potenziometro);
%\item Breadboard e cablaggi vari.
\end{itemize}

\subsection{Misura dello Slew Rate}

\begin{wrapfigure}[15]{r}{0.55\textwidth}
  \begin{center}
    \includegraphics[width=0.32\textwidth]{../E03/latex/slew_rate.pdf}
  \end{center}
  \caption{Schema del circuito utilizzato per stimare lo Slew Rate. La resistenza utilizzata è $R=\SI{2.17\pm0.01}{\kilo\ohm}$; la capacità $C=\SI{200 \pm 10}{\pF}$.}
  \label{cir3:slew_rate}
\end{wrapfigure}

Lo Slew Rate è uno dei parametri che caratterizzano l'amplificatore operazionale ed è definito come
$$SR = \mathrm{max}\left[\frac{\Delta V}{\Delta t}\right]$$
Esso indica la massima velocità con cui l'opamp varia la tensione in uscita nell'unità di tempo.
Ciò significa che, per segnali in entrata con derivata temporale superiore ad SR dell'amplificatore operazionale, avremo un segnale distorto in uscita.

Per valutare la risposta in uscita del nostro $\mu$A741 abbiamo utilizzato il circuito in Figura \ref{cir3:slew_rate} fornendo in ingresso un'onda quadra caratterizzata da una velocità di variazione dei segnale di $SR_{gen}=\SI{315}{\volt\per\micro\second}$ (misurata tramite l'oscilloscopio).
Essendo $SR_{gen}$ diversi ordini di grandezza maggiore del valore atteso di $SR$ ($\approx \SI{0.5}{\volt\per\micro\second} $) possiamo considerare l'onda quadra come ideale, cioè avente pendenza infinita.
La capacità posta in parallelo alla resistenza è utilizzata come serbatoio di cariche per attenuare fenomeni di rumore nel segnale creati dal generatore.

\begin{figure}[ht]
 \centering
   {\includegraphics[width=\textwidth]{../E03/latex/sr_uad.pdf}}
 \caption{Grafico della tensione in uscita (rosso) ed in entrata (nero) in funzione del tempo. I $\Delta$ sono relativi alle differenze fra le coordinate dei punti segnati. Si noti che le scale dei tempi sono differenti per i due grafici.}
 \label{gr3:slew_rate}
\end{figure}

Abbiamo posto inizialmente la frequenza a $f=1$\si{\kilo\hertz} e la tensione picco-picco a $V_{pp}=10$ \si{\volt};
abbiamo poi misurato con i cursori dell'oscilloscopio la differenza fra il 10\% e il 90\% del valore di tensione picco picco in uscita $\Delta V = (8.07 \pm 0.01)$ \si{\volt}  e la relativa $\Delta t = (15.9 \pm 0.1)$ \si{\micro\second}, ottenendo uno $SR= \frac{\Delta V}{\Delta t} = \SI{0.507 \pm 0.003}{\volt\per\micro\second}$.

Però, come si può notare dal grafico in Figura \ref{gr3:slew_rate}, nella zona vicina al 10\% è presente una parte di segnale affetta da rumore che potrebbe rendere poco precisa la nostra misura.
Per ovviare a questo problema abbiamo spostato il limite inferiore ad un valore percentuale più alto (25\%).
Quest'operazione è permessa dal fatto che siamo interessati al rapporto fra i $\Delta$ e questo rimane costante su una retta.
Portando dunque la percentuale del limite inferiore al 25\% abbiamo ottenuto i seguenti valori: $\Delta V = \SI{6.48 \pm 0.01}{\volt}$ e $\Delta t = \SI{14.5 \pm 0.1}{\ms}$, dunque $SR = \SI{0.447 \pm 0.003}{\volt\per\ms}$. Durante l'esperienza utilizzeremo questo valore come riferimento per lo slew rate del nostro amplificatore, in quanto più restrittivo.

Abbiamo anche notato che lo slew rate è differente a seconda che il segnale abbia derivata $dV/dt$ positiva o negativa (grafico in Figura \ref{gr3:slew_rate}, immagine di destra).
Nel secondo caso, infatti, abbiamo ottenuto i seguenti valori: $\Delta V = (8.00 \pm 0.01)$ \si{\volt} e $\Delta t = (23.2 \pm 0.1)$ \si{\micro\second}, quindi $SR = (0.345 \pm 0.004)$ \si{\volt\per\micro\second}.

\subsection{Misura della capacità di pilotaggio}

\begin{wrapfigure}[18]{l}{0.55\textwidth}
  \begin{center}
    \includegraphics[width=0.33\textwidth]{../E03/latex/max_current.pdf}
  \end{center}
  \caption{Schema del circuito utilizzato per misurare la corrente massima. La resistenza utilizzata è $R=\SI{98.47\pm0.01}{\ohm}$; il generatore fornisce un'onda triangolare di $f=1$ \si{\kilo\hertz} e $5 V_{pp}$.}
  \label{cir3:max_current}
\end{wrapfigure}

Un altro limite dell'amplificatore operazionale è quello di avere un valore massimo di corrente erogabile in uscita.
Pertanto, utilizzando un circuito di feedback come in Figura \ref{cir3:max_current}, la corrente massima e il valore della resistenza $R$ fisseranno la tensione letta ai capi della resistenza stessa, secondo questa relazione:
\begin{equation*}
	I_{max} = \frac{V_{clip}}{R}
\end{equation*}
dove $V_{clip}$ è la tensione in cui il segnale in uscita risulta tagliato rispetto a quello in ingresso (fenomeno del \textit{clipping}).

Invertendo la precedente equazione, abbiamo dunque potuto misurare $I_{max}$, ponendo una resistenza di carico fra l'uscita dell'operazionale e terra\footnote{Si noti che, data l'alta impedenza in ingresso dell'oscilloscopio (\SI{1}{\mega\ohm}), non era possibile misurare la corrente massima con questo strumento (attesa, come da specifiche del costruttore, sui \SI{15}{\milli\ampere}).}, con l'oscilloscopio ai capi di tale resistenza per la misura di tensione (circuito in Figura \ref{cir3:max_current}).

\begin{figure}[ht]
 \centering
   {\includegraphics[width=\textwidth]{../E03/latex/clip.pdf}}
 \caption{Grafico della tensione in uscita (verde) ed in entrata (nero) in funzione del tempo. Si nota subito che il clip a tensioni negative avviene a tensioni in modulo maggiori.}
 \label{gr3:clip}
\end{figure}

Bisogna però considerare che, essendo il segnale in entrata in tensione alternata, la corrente risulta oscillare fra valori negativi e positivi (corrente rispettivamente entrante o uscente dal nodo A verso $R$.
Si può notare anche solo osservando l'asimmetria nel grafico in Figura \ref{gr3:clip} che i valori per le due correnti sono in modulo differenti.
Se consideriamo $V^+ = \SI{1.92 \pm 0.01}{\volt}$ la tensione di clip positiva e $V^- = \SI{3.45 \pm 0.01}{\volt}$ quella negativa, otteniamo che le correnti massime erogabili dall'opamp sono:

$$I_{V^+} = \frac{V^+}{R} = \SI{19.5 \pm 0.1}{\mA}  \qquad I_{V^-} = \frac{V^-}{R} = \SI{35.2 \pm 0.2}{\mA}$$
entrambi i valori sono compatibili con quelli del costruttore (valori di corrente fino a $\pm 40$ \si{\milli\volt}).

L'asimmetria delle correnti può essere spiegata considerando l'asimmetria della circuiteria interna dell'opamp, la quale è composta da diversi transistor che solo idealmente possono essere tutti delle medesime ed identiche caratteristiche.

\newpage
\subsection{Verifica della banda passante}
\label{par3:bode}

In questa parte dell'esperienza vogliamo valutare la banda passante di un circuito amplificatore (non invertente) considerando che il guadagno a maglia aperta dell'amplificatore operazionale dipende della frequenza.

\subsubsection*{Funzione di trasferimento e frequenza di taglio}

Esaminiamo dapprima il circuito, calcolandone la funzione di trasferimento. Essa è definita come:\phantom{xxxxxxxxxxxxxxxxx}
\begin{wrapfigure}[20]{l}{0.45\textwidth}
  \begin{center}
    \includegraphics[width=0.32\textwidth]{../E03/latex/bandwidth.pdf}
  \end{center}
  \caption{Circuito utilizzato per valutare la banda passante. $R_1=\SI{997\pm1}{\ohm}$ è fissa; mentre $R_2$ è stata cambiata con valori $R_2^{11\mathrm{x}}=\SI{9.91 \pm 0.01}{\kilo\ohm}$ e $R_2^{101\mathrm{x}}=\SI{99.1 \pm 0.1}{\kilo\ohm}$ a seconda del guadagno desiderato.}
  \label{cir3:banda}
\end{wrapfigure}
\begin{equation}
H(s)=\frac{V_{out}}{V_{in}}
\label{eq3:def_f_trasf}
\end{equation}
dove $s=j\omega$.
Per ogni amplificatore operazionale vale
\begin{equation}
V_{out}=A(s) (V^+-V^-)
\label{eq3:regola_opamp}
\end{equation}
con il guadagno che dipende da $s=j\omega$.
Consideriamo
$$V^+ = V_{in} \qquad V^-=V_{out} \frac{R_1}{R_1+R_2}$$
dove $V^-$ è ricavato dall'analisi delle correnti sul ramo di retroazione.
Sostituendo questi valori in (\ref{eq3:regola_opamp}) otteniamo
$$A(s) V_{in} = V_{out} + V_{out} A(s) \beta \quad con \quad \beta = \frac{R_1}{R_1+R_2} = \frac{1}{G}$$
Si noti che $G$ è il guadagno circuitale di un amplificatore non invertente.
Si può ora calcolare $H(s)$:
\begin{equation}
H(s)=\frac{V_{out}}{V_{in}}=\frac{1}{\beta}\frac{1}{1+\frac{1}{A(s) \beta}}
\label{eq3:funz_trasfe}
\end{equation}
da cui è facile notare che per $A(s) \rightarrow + \infty$ si ottiene l'approssimazione di amplificatore ideale $H(s)=G=\frac{1}{\beta}=1+\frac{R_2}{R_1}$ e l'equazione diventa indipendente dal guadagno a maglia aperta $A$.

Se schematizziamo l'opamp come un filtro passa basso, abbiamo che il guadagno a maglia aperta varia con la frequenza secondo la legge
\begin{equation}
A(j\omega)=\frac{A_{ol}}{1+j\frac{\omega}{\omega_0}}
\label{eq3:passa_basso}
\end{equation}
con $A_{ol}$ il guadagno dell'operazionale per segnali costanti ($f = 0$) e $\omega_0 \approx 8$ \si{\hertz} la prima frequenza di taglio dell'operazionale data dalla capacità di compensazione nella circuiteria interna. Sostituendo questa relazione in (\ref{eq3:funz_trasfe}), abbiamo
\begin{equation}
H(j\omega)=\frac{\frac{A_{ol}}{1+A_{ol}\beta}}{1+j \frac{\omega}{(1+A_{ol}\beta)\omega_0}} \approx \frac{G}{1+j\frac{\omega}{\omega_t}}
\label{eq3:funzione_trasferimento}
\end{equation}
da cui si nota che la nuova frequenza di taglio è data da
\begin{equation}
\omega_t=(1+A_{ol}\beta)\omega_0
\label{eq3:freq_taglio}
\end{equation}
Nel circuito in Figura \ref{cir3:banda} abbiamo impostato un guadagno $G$ retroazionato di $11$ e $101$, quindi ci aspettiamo\footnote{Questi valori sono approssimati perché cercheremo di calcolare $A_{ol}$ con le frequenze a nostra disposizione, ricavate con l'analisi della banda, piuttosto che calcolare le frequenze di taglio supponendo certo $A_{ol}$.} frequenze di taglio, ricavate con l'equazione precedente, rispettivamente sull'ordine dei $10^5$ e $10^4$.

Un secondo modo per fare una stima delle frequenze di taglio è quello di sfruttare la costante \textit{gain-bandwidth product} GBWP $= \omega_t G$.

\subsubsection*{Effetto dello Slew Rate}

Durante la misurazione bisogna però tenere conto dello slew rate.
Esso infatti può influenzare negativamente la misura distorcendo la forma d'onda in uscita.
È necessario pertanto scegliere valori di tensione del segnale in ingresso appropriati.
Supponiamo dunque di avere una forma d'onda in uscita del tipo $V_{out} \sin (2 \pi f_0 t) = V_{in} G \sin (2 \pi f_0 t)$.
Affinché non si raggiunga lo SR, deve essere verificato
$$\mathrm{max}\left[ \frac{dV_{out}}{dt} \right] = \mathrm{max} \left[ 2 \pi V_{in} G f_0 \cos (2 \pi f_0 t) \right]= 2 \pi V_{in} G f_0 < 0.35 \si{\volt\per\micro\second}$$
dove l'ultimo valore è stato scelto per tenerci sempre abbastanza distanti dal valore reale di SR. Dunque vale che
\begin{equation}
V_{in}<\frac{0.35 \si{\volt\per\micro\second}}{2 \pi G f_0} = \frac{0.35 \si{\volt\per\micro\second}}{2 \pi \mathrm{GBWP}}
\label{eq3:tensione_max}
\end{equation}
Considerando $f_0$ come la frequenza di taglio stimata prima, otteniamo che i valori di tensione da usare nei due casi è di circa $V_{in}^{11\mathrm{x}} = 100$ \si{\milli\volt}$= V_{in}^{101\mathrm{x}}$. Si nota che la regola \ref{eq3:tensione_max} è valida quando il guadagno non dipende dalla frequenza: infatti dopo la frequenza di taglio il guadagno del circuito $G$ cala e dunque i valori di tensione ammessi risultano di valore maggiore.

\subsubsection*{Grafici di Bode}
\label{par3:sub_bode}

\begin{figure}[ht]
 \centering
   {\includegraphics[width=0.75\textwidth]{../E03/latex/bode.pdf}}
 \caption{Diagrammi di Bode per il guadagno del circuito in Figura \ref{cir3:banda}. Le leggi plottate sono (\ref{eq3:fit_gain}) e (\ref{eq3:fit_fase}), con i rispettivi guadagni $G$ del circuito.}
 \label{gr3:bode}
\end{figure}

Utilizzando una forma d'onda sinusoidale abbiamo dunque misurato la tensione di uscita del circuito e la fase fra i segnali, per poi plottare i grafici di Bode (Figura \ref{gr3:bode}).

Dalla funzione di trasferimento (\ref{eq3:funzione_trasferimento}) possiamo inoltre trovare le relazioni fra il guadagno e la fase con la frequenza. Sappiamo infatti che il guadagno è dato $V_{out}/V_{in}$, ed in generale queste due tensioni possono essere considerate come complesse. Se siamo dunque interessati al guadagno dobbiamo considerare piuttosto il rapporto dei moduli, ed otteniamo la seguente relazione
\begin{equation}
\mathrm{Gain}=\frac{|V_{out}|}{|V_{in}|}=|H(j\omega)|=\sqrt{\frac{A_{ol}^2 \omega_0^2}{(\omega_0 + A_{ol} G \omega_0)^2 + \omega^2}}
\label{eq3:fit_gain}
\end{equation}
mentre, per la fase utilizziamo la solita formula della fase per un numero complesso (nel nostro caso $V_{out}/V_{in}$), ed otteniamo
\begin{equation}
\varphi=\arctan\left[{\frac{\mathrm{Im}[H(j\omega)]}{\mathrm{Re}[H(j\omega)]}}\right]= - \arctan\left[\frac{\omega}{(1+A_{ol}G)\omega_0}\right]
\label{eq3:fit_fase}
\end{equation}

Date queste leggi è possibile stimare $A_{ol}$ e $\omega_0$ dell'operazionale: rimandiamo al paragrafo \ref{par3:A_ol} per questi calcoli, dove saranno confrontati i valori ottenuti con questi dati e quelli presi al paragrafo \ref{par3:open}.

Infine, i valori di frequenza di taglio possono essere calcolati dalla legge teorica nel punto di flesso ($45^{\circ}$). Questi sono: $f_t^{101\mathrm{x}} = (10.3 \pm 0.1)$ \si{\kHz} e $f_t^{11\mathrm{x}} = (92 \pm 1)$ \si{\kHz}, in accordo con gli ordini di grandezza stimati in precedenza.

\newpage
\subsection{Guadagno Open Loop}
\label{par3:open}

\begin{wrapfigure}[17]{r}{0.45\textwidth}
  \begin{center}
    \includegraphics[width=0.280\textwidth]{../E03/latex/LF_ol.pdf}
  \end{center}
  \caption{Circuito utilizzato per misurare il guadagno $A_{ol}$. I valori di $R_4$ utilizzati sono $\SI{5.2\pm0.1}{\ohm}$,  $\SI{9.6\pm0.1}{\ohm}$, $\SI{101.3\pm0.1}{\ohm}$ e $\SI{998.2\pm0.1}{\ohm}$ mentre $R_1=\SI{98\pm1}{\kohm}$, $R_2=\SI{97\pm1}{\kohm}$ e $R_3=\SI{98\pm1}{\kohm}$.}
  \label{cir3:low_frequency}
\end{wrapfigure}

Come ultima parte dell'esperienza, abbiamo cercato di misurare il guadagno \textit{open-loop} $A_{ol}$ del nostro amplificatore reale.
Come sappiamo, tale guadagno è una funzione della frequenza e dobbiamo dunque trovare un modo per misurarlo.

L'opamp $\mu$A741 è progettato per avere un guadagno costante di circa \num{2E5} fino alla frequenza di circa \SI{8}{\Hz}.
Oltre tale frequenza abbiamo una caduta di circa \SI{8}{\decibel}/ottava (\SI{20}{\decibel}/decade) ed è pertanto impossibile effettuare misure a loop aperto in quanto piccolissime variazioni di tensione ai due ingressi causerebbero grandi effetti in uscita, mandando l'integrato in saturazione.

Per questo motivo abbiamo utilizzato il circuito riportato in Figura \ref{cir3:low_frequency}, in modo da controllare la tensione in uscita e limitarla a valori inferiori alla tensione di saturazione.
Utilizzando l'oscilloscopio abbiamo poi misurato i valori di tensione $V_A$ e $V_{out}$.

Cerchiamo dunque di legare il guadagno a tali quantità.
Nel caso dell'amplificatore operazionale sappiamo che vale l'eq. (\ref{eq3:regola_opamp}), dove $(V_+-V_-)$ è semplicemente la differenza di potenziale presente tra i due ingressi, che sarà ovviamente data da $\Delta V = I_3R_4$, dove $I_3$ è la corrente che scorre attraverso $R_3$ ed $R_4$\footnote{Assumiamo che la corrente assorbita dagli ingressi sia trascurabile.}.
È quindi facile stimare $I_3=\frac{V_A}{R_3+R_4}$ conoscendo $V_A$.

Si ottiene dunque:
\begin{equation}
A_{ol}=\frac{{V_{out}}^{pp}}{{V_A}^{pp}} \frac{R_4+R_3}{R_4}
\label{eq3:lfgain}
\end{equation}

Osserviamo come in equazione (\ref{eq3:lfgain}) non compaia il termine $V_{in}$.
Dovremo dunque scegliere per le varie frequenze una tensione picco-picco in ingresso adeguata in modo che il nostro opamp non saturi e allo stesso tempo $V_A$ sia sufficientemente grande da essere poco influenzata dal rumore di fondo e sia stabile il più possibile.

Durante le misurazioni abbiamo deciso di variare la resistenza $R_4$ suddividendo il procedimento di misura in diversi step; questa strategia ci ha permesso di mantenere i valori di $V_A$ e $V_{out}$ sempre nell'ordine del volt.
I dati ricavati sono esposti nel grafico in Figura \ref{cir3:gain_open_loop}.

\subsubsection*{Misura per alte frequenze}

\begin{wrapfigure}[13]{r}{0.45\textwidth}
  \begin{center}
    \includegraphics[width=0.28\textwidth]{../E03/latex/HF_ol.pdf}
  \end{center}
  \caption{Circuito utilizzato per stimare il guadagno $A_{ol}$ a basse frequenze. I valori delle resistenze utilizzate sono le stesse del circuito in Figura \ref{cir3:low_frequency}.}
  \label{cir3:high_frequency}
\end{wrapfigure}

Per non dover utilizzare resistenze $R_4$ troppo grandi ad alte frequenze ($>$ \SI{70}{\kilo\hertz}), abbiamo provato ad effettuare le misure direttamente ad open-loop, dato che il guadagno è contenuto (\num{<100}).
Tuttavia il generatore di forme d'onda non emette un segnale pulito alla sola frequenza voluta, in quanto anch'esso è uno strumento reale e dunque affetto da errori.
Esso potrebbe dunque emettere segnali a frequenze basse o addirittura un piccolo offset.
Poiché per tali frequenze il guadagno dell'amplificatore operazionale è enorme, il segnale in uscita viene immediatamente distorto da tali rumori indesiderati.
Non è da sottovalutare poi la presenza dei residui dell'offset dell'opamp.
Esso infatti dipende anche dalla temperatura e dunque la sua compensazione non potrà mai essere perfetta ad ogni temperatura.
Misurare il guadagno $A_{ol}$ effettivamente a open-loop è particolarmente problematico.

Per aggirare tali ostacoli possiamo utilizzare il circuito riportato in Figura \ref{cir3:low_frequency}, rimuovendo la resistenza $R_4$ e sostituendo $R_3$ con un filo partendo dalla configurazione di Figura \ref{cir3:high_frequency}.
Cosi facendo otteniamo la semplice relazione $A_{ol}=\frac{{V_{out}}^{pp}}{{V_A}^{pp}}$, che possiamo anche vedere come limite di eq. (\ref{eq3:lfgain}) per $R_4 \rightarrow \infty$.
I dati ottenuti sono riportati nel grafico in Figura \ref{cir3:gain_open_loop}.

\begin{figure}[h]
	\centering
	\includegraphics[width=0.9\textwidth]{../E03/latex/gol.pdf}
	\caption{Grafico semilogaritmico del guadagno (in $dB$) in funzione della frequenza. Ad ogni forma dei punti (e relativo colore) corrisponde una diversa resistenza $R_4$ utilizzata. Per quanto riguarda $R_4=\infty$, intendiamo la mancanza del collegamento fra ingresso invertente e non invertente (circuito in Figura \ref{cir3:high_frequency}). La legge plottata è la (\ref{eq3:fit_gain_ol}).}
  \label{cir3:gain_open_loop}
\end{figure}

Anche in questo caso, considerando l'operazionale come un semplice passa basso, possiamo utilizzare l'equazione (\ref{eq3:passa_basso}) e trovare la relazione fra il guadagno e la frequenza. Otteniamo
\begin{equation}
\mathrm{Gain}=\frac{|V_{out}|}{|V_{in}|}=\sqrt{\frac{A_{ol}^2 \omega_0^2}{\omega^2 + \omega_0^2}}
\label{eq3:fit_gain_ol}
\end{equation}

Per un fit e i parametri si veda il paragrafo successivo.

\newpage
\subsection{Calcolo di $A_{ol}$}
\label{par3:A_ol}

\subsubsection{$A_{ol}$ da circuiti retroazionati}

Abbiamo ottenuto una stima di $A_{ol}$ e $\omega_0$ utilizzandoli come parametri su un fit delle leggi (\ref{eq3:fit_gain}) e (\ref{eq3:fit_fase}) sui dati sperimentali.
Inoltre, dato che l'opamp utilizzato era lo stesso per tutte le misurazioni, con una media pesata sui valori ottenuti con i quattro fit (due con (\ref{eq3:fit_gain}) e due con (\ref{eq3:fit_fase})), abbiamo ottenuto i seguenti valori sperimentali
$$A_{ol}=(1.8\pm0.1)\times 10^5 \qquad f_0=(5.9\pm0.2)\si{\hertz}$$

Ovviamente possiamo ottenere un valore di $A_{ol}$ anche utilizzando la (\ref{eq3:freq_taglio})
$$A_{ol} = \frac{G}{\omega_0} (\omega_t - \omega_0)$$
e considerando la frequenza $f_0$ ricavata dal fit. Si ottiene un valore teorico (media dei valori con i due guadagni del circuito) di $A_{ol} = (1.71\pm 0.06)\times 10^5$, compatibile con il valore sperimentale.

\subsubsection{$A_{ol}$ da open loop}

Con il fit dalla (\ref{eq3:fit_gain_ol}) otteniamo i seguenti valori
$$A_{ol}=(2.0\pm0.1)\times 10^5 \qquad f_0=(5.8\pm0.2)\si{\hertz}$$

I dati ottenuti non risultano compatibili, riteniamo che ciò sia dovuto ad errori parassiti non ben identificati.

\subsection*{Conclusioni}

In questa esperienza siamo riusciti a stimare lo "slew rate" di un amplificatore operazionale $\mu$A741 e la massima corrente erogata. Tali valori sono risultati compatibili con quelli riportati sul datasheet del costruttore. 

Successivamente abbiamo verificato la banda passante per un circuito non invertente retroazionato negativamente con guadagno G=11 e G=101. Dai diagrammi di Bode abbiamo visto come aumentando il Gain diminuiva la banda passante (ovvero GBWP=const).

Infine abbiamo valutato il guadagno open-loop a diverse frequenze, utilizzando dei circuiti che ci permettevano un controllo del guadagno effettivo e dunque di evitare problemi quali la saturazione dell'amplificatore operazionale. Il guadagno open-loop per tensioni costanti è stato stimato essere un valore compreso fra \num{1.8E5} e \num{2.0E5}, valore compatibile con quelli forniti dal costruttore.
